% DTW教学算例的累计代价表；宽约85mm。x=(0,0,2,4,4)，y=(0,2,4)，局部代价|x_i-y_j|。
\begin{tikzpicture}[x=1.25cm,y=1.0cm,font=\figurefont\small,
  arr/.style={-{Stealth[length=2mm]},BookTeal,line width=1.0pt}]
  % 网格：列为视频帧i=1..5，行为音频单元j=1..3
  \foreach \i/\xv in {1/0,2/0,3/2,4/4,5/4}{
    \node at (\i,-0.05) {$x_{\i}=\xv$};
  }
  \foreach \j/\yv in {1/0,2/2,3/4}{
    \node[anchor=east] at (0.4,\j) {$y_{\j}=\yv$};
  }
  \foreach \i/\j in {1/1,2/1,3/2,4/3,5/3}{
    \fill[BookTeal!15] (\i-0.5,\j-0.45) rectangle (\i+0.5,\j+0.45);
  }
  \foreach \i in {1,...,5}{
    \foreach \j in {1,2,3}{
      \draw[BookRule,line width=0.5pt] (\i-0.5,\j-0.45) rectangle (\i+0.5,\j+0.45);
    }
  }
  \foreach \i/\j/\v in {1/1/0,1/2/2,1/3/6,
                        2/1/0,2/2/2,2/3/6,
                        3/1/2,3/2/0,3/3/2,
                        4/1/6,4/2/2,4/3/0,
                        5/1/10,5/2/4,5/3/0}{
    \node at (\i,\j) {$\v$};
  }
  % 回溯得到的规整路径
  \draw[arr] (1.2,1) -- (1.8,1);
  \draw[arr] (2.2,1.15) -- (2.8,1.85);
  \draw[arr] (3.2,2.15) -- (3.8,2.85);
  \draw[arr] (4.2,3) -- (4.8,3);
  \node[anchor=north] at (3,-0.4) {视频帧动作强度$x_i$};
  \node[rotate=90,anchor=south] at (-0.75,2) {音频能量$y_j$};
\end{tikzpicture}
